\documentclass[11pt,a4paper]{article}
\usepackage[margin=1in]{geometry}
\usepackage[T1]{fontenc}
\usepackage[utf8]{inputenc}
\usepackage{lmodern}
\usepackage{microtype}
\usepackage{enumitem}
\usepackage{hyperref}
\usepackage{fancyhdr}
\usepackage{array}
\usepackage[authoryear,round]{natbib}
\setlength{\parindent}{0pt}
\setlength{\parskip}{0.65em}
\pagestyle{fancy}
\fancyhf{}
\fancyfoot[C]{\thepage}
\renewcommand{\headrulewidth}{0pt}
\newcommand{\guidance}[1]{\textit{#1}}
\begin{document}
\begin{center}
{\LARGE\bfseries AI Collusion Forensics: Comparing Communication and Behavioral Evidence}\par
\vspace{1.5em}
\begin{tabular}{>{\centering\arraybackslash}p{0.28\textwidth} >{\centering\arraybackslash}p{0.28\textwidth} >{\centering\arraybackslash}p{0.28\textwidth}}
Agnes Ahalya Arogyaraj \\
Saarland University \\[1.3em]

\end{tabular}
\par\vspace{1em}
With\\Apart Research
\end{center}
\begin{abstract}
\guidance{Summarize your project in 150--250 words. A strong abstract lets a reviewer understand what you did and why it matters without reading anything else. Make sure to cover: the problem, your approach, key results, and the main takeaway. Polish it last: the abstract should reflect your final results, not your initial plan.}
\end{abstract}
\vspace{1em}
\noindent\textbf{How to use this template:} Replace the italicized guidance text under each section with your content. The section structure is strong guidance but not rigid. If your project requires a different organization, feel free to adapt. Delete all guidance text including this info box before submitting. Make sure the project title and author information above are up to date.

\noindent\textbf{How your submission will be evaluated:} Your project will be judged on the quality of this written report. To understand what we look for, see our \href{https://apartresearch.com/}{Evaluation Rubric}.

\noindent\textbf{Recommended length:} 4 pages excluding references and appendix.\\
\textbf{Rough guide:} Intro \& Related Work 1p, Methods and Results: 2.5p, Discussion 0.5p.

\vfill
\noindent\textsuperscript{1} Research conducted at the AI Collusion Research Sprint, October 2026
\newpage

\section{Introduction}
\guidance{What problem are you addressing and why does it matter?\\
Connect it to your work: we want to know why your work is practically valuable.}

\guidance{Provide enough background for readers to understand your work.\\
If relevant, briefly describe the threat model or failure mode you're addressing; reference prior work that motivates why it is worth addressing, or explain it yourself.}

\guidance{Aspire to clearly list your most important contributions that go beyond what exists today.}

Our main contributions are:
\begin{enumerate}
\item \guidance{[First contribution --- what new thing did you create, discover, or demonstrate?]}
\item \guidance{[Second contribution]}
\item \guidance{[Third contribution, if applicable]}
\end{enumerate}

\section{Related Work}
Research on multi-agent systems has examined how communication and behavior reveal coordination. COLOSSEUM(\cite{nakamura2026colosseum}) investigates collusion through agent messages and observed actions, distinguishing expressed intentions from executed behavior . \cite{radev2026best} show that access to agent communication can improve detection of coordinated sabotage. \cite{shi2026agents} examine differences between private intentions, public commitments, and final actions in repeated games, providing the experimental foundation for our study.

Monitoring research also highlights how evidence presentation affects detection. \cite{arike2026does} find that additional information does not always improve sabotage detection, while \cite{nguyen2026traceguard} introduce TraceGuard, which uses structured assessments of agent behavior. Building on these findings, we compare four monitoring approaches on identical multi-agent trajectories: (1) public messages only, (2) actions and outcomes only, (3) combined messages and actions, and (4) structured combined evidence comparing stated and observed behavior. We investigate whether combining communication and behavioral evidence improves collusion detection over either source alone.

\section{Methods}
\guidance{Describe your approach clearly enough that someone could replicate it. Include key design choices and justify them where relevant (hint: the more you can back up your design choices by referencing prior work, the better).}

\guidance{What models/datasets/tools did you use and why? What were key parameters or design decisions? What did you try that didn't work? Could someone reproduce your work from this description?}

\section{Results}
\guidance{Present your main findings with appropriate evidence. Use figures and tables where appropriate (we strongly encourage at least one figure, see tips below). Distinguish between observations and interpretations.}

\guidance{Argue why your claims are robust. E.g., if your approach ``performs better'' than alternatives:}
\begin{enumerate}[label=(\alph*)]
\item \guidance{Do you have enough data? Is the difference statistically significant?}
\item \guidance{Is it robust? Or do small changes to your setup cause substantial changes to your results?}
\end{enumerate}

\begin{quote}
\textbf{Tips for figures and tables:}
\begin{itemize}[leftmargin=*]
\item Number all figures (Figure 1, Figure 2...) and tables (Table 1, Table 2...)
\item Include descriptive captions that can be understood without the main text
\item Place figures/tables near where they're first referenced
\item Ensure text in figures is legible!
\end{itemize}
\end{quote}

\section{Discussion and Limitations}
\guidance{Discuss the broader implications for AI safety.\\
What do your results mean? What trends do you notice and what might they indicate?}

\subsection*{Limitations}
\guidance{What are the limitations of your work? What threat models or failure modes did you not address? Be honest about constraints --- methodological limitations, scope limitations, or aspects you couldn't fully address in the hackathon timeframe. Explicitly note the assumptions you made, whether implicitly or explicitly, and how the interpretation of your results would change if a given assumption did not hold.}

\subsection*{Future Work}
\guidance{What are the natural next steps? How could this work be extended?}

\section{Conclusion}
\guidance{Briefly summarize your main findings and their implications (1--2 paragraphs).}

\subsection*{Code and Data}
\guidance{Include links if applicable. If your project doesn't involve code (e.g., policy analysis) or if there are info-hazard considerations, note that here.}
\begin{itemize}
\item Code repository: [Link to GitHub/GitLab if applicable]
\item Data/Datasets: [Link if applicable]
\item Other artifacts (optional): [Demo link, video walkthrough, Hugging Face Space, etc.]
\end{itemize}



\bibliographystyle{plainnat}
\bibliography{references} 



\section*{LLM Usage Statement}
Large language models (LLMs) were used to brainstorm research ideas, summarize and understand related literature, assist with code generation and debugging, and support drafting and editing the report. The research design, implementation decisions, interpretation of results, and final conclusions remained the responsibility of the author. All AI-generated content was reviewed before inclusion.

\end{document}
